\chapter{Scripting the metadata itself}
\label{ch:scriptable}

The \lang{dynamic} backend makes your library's metadata available \emph{to
C++}. Every consumer of it in the repository is C++: the terminal interpreter,
the Qt viewer, the ImGui runtime, the QML reflected object --- hand-written
walkers, one per toolkit.

The last step is to bind \emph{that API} through rosetta's own backends, so the
walker can be written in Python, Lua or JavaScript instead. This is the
Graphite/GOM move: expose the object model itself to the scripting layer, and
the user interface becomes a script rather than a backend.

\section{The trick, in one sentence}

Run the pipeline twice, and make the first run's \emph{output} the second run's
\emph{input library}.

\begin{center}
\begin{tikzpicture}[
  node distance=5mm,
  b/.style={draw,rounded corners=2pt,align=center,font=\scriptsize,inner sep=4pt,text width=30mm,minimum height=9mm},
  user/.style={b,fill=blue!6,draw=rosettablue},
  gen/.style={b,fill=black!4,draw=black!35,dashed},
  fin/.style={b,fill=green!7,draw=green!45!black},
  arr/.style={-{Stealth[length=2mm]},draw=black!55,font=\scriptsize}
]
\node[user] (h)  {\texttt{scene.h}\\\scriptsize your library};
\node[gen,right=10mm of h] (d) {\texttt{auto\_dynamic.\{h,cpp\}}\\\scriptsize your types as DATA};
\node[user,below=9mm of h] (s) {\texttt{<rosetta/script.h>}\\\scriptsize the fixed facade};
\node[fin,right=10mm of s] (m) {one module per language\\\scriptsize names no bound type};
\draw[arr] (h) -- node[above]{\ \ stage 1} (d);
\draw[arr] (s) -- node[above]{\ \ stage 2} (m);
\draw[arr] (d) -- node[right]{\texttt{user\_sources}} (m);
\end{tikzpicture}
\end{center}

By stage~2 your types are \textbf{data}, not C++. So stage~2's manifest never
names one, and the resulting module works for \emph{any} library that has been
through stage~1.

\section{Why a facade rather than binding \code{dynamic.h}}

\code{rosetta::dyn} cannot be fed to the backends as it stands. Three shapes
block it:

\begin{center}
\small
\begin{tabular}{@{}L{3.9cm}L{4.6cm}L{4.5cm}@{}}
\toprule
\textbf{In \code{dynamic.h}} & \textbf{Why it blocks} & \textbf{In
  \code{script.h}} \\
\midrule
\code{MetaField *fields; size\_t n\_fields} &
  pointer-and-count, so vector marshalling never fires &
  \code{std::vector<FieldInfo> fields()} \\
\addlinespace
\code{const MetaClass *}, \code{const TypeDesc *} &
  raw pointer to an unbound type --- the gate skips it &
  value-semantics handles \\
\addlinespace
\code{using Thunk = Any (*)(\ldots)} &
  function pointers marshal nowhere &
  ordinary \code{get} / \code{set} / \code{call} \\
\bottomrule
\end{tabular}
\end{center}

\code{<rosetta/script.h>} is that reshaping and nothing more. Each handle wraps
\textbf{one \code{const} pointer} into the tables, so binding it copies no
metadata --- the \code{.rodata} guarantee survives.

\begin{note}
Nothing in the facade is named \code{Class}, \code{Type}, \code{Object} or
\code{Method}, because those collide in the target languages
(\code{java.lang.Class}, JavaScript's \code{Object}, Lua's \code{type}). Hence
\code{ClassInfo}, \code{TypeInfo}, \code{Instance}, \code{MethodInfo}.
\end{note}

Errors come back as \code{Outcome} (an ok flag plus text), never as an
exception. Unwinding through a foreign VM's frames is how you get a crash
instead of a \code{TypeError}.

\section{The one hand-written piece}

\code{Value} is type erasure --- one C++ type standing for ``whatever the host
language just handed us''. Reflection can \emph{describe} it but cannot know
that a Python \code{float}, a Lua number and a JavaScript \code{Number} should
all become the same box. That mapping is per-language knowledge, so it is
written per language, in
\code{runtime/script/\{python,lua,node,wasm\}.h}.

Only half of each caster is per-language, though:

\begin{itemize}
  \item \textbf{Reading} a host value is irreducible --- only Python knows what
    \code{PyFloat\_Check} means.
  \item \textbf{Writing} one is the same shape everywhere, so it lives once in
    \code{rosetta::script::visit()}, and each caster supplies a sink of nine
    small methods (\code{on\_none}, \code{on\_bool}, \code{on\_int},
    \code{on\_real}, \code{on\_string}, \code{on\_enum}, \code{on\_list},
    \code{on\_object}, \code{on\_unknown}).
\end{itemize}

No two casters can drift apart on whether an enum crosses as its integer, or on
what an unclaimed type renders as.

The casters reach the generated code through \field{module\_init}'s
\field{headers}, which is the only hook that emits an \code{\#include}
\emph{after} the framework header and \emph{before} the bound classes --- exactly
the window a type caster needs. Because that list is global rather than
per-target, each file guards its whole body on a macro only its own target
defines.

\begin{warn}
One subtlety the type gate forces: \code{Value} must stay in the bound
\field{classes} list. Bindability is decided at generation time and cannot know
a caster will exist, so dropping \code{Value} would silently skip every member
that mentions it --- \code{Outcome::value}, \code{Instance::set},
\code{Instance::call}, all of it. The cost is one vestigial \code{Value} type
per module that nothing ever produces.
\end{warn}

\section{What a project has to write}

Nothing, beyond where its own files are. The binding surface is fixed, so it
ships as manifest data in \code{include/rosetta/presets/scriptable.json} and a
manifest names it:

\begin{lstlisting}[style=json]
{
    "preset": "scriptable",
    "rosetta_include": "../../include",
    "user_include": ["../lib", "../lib/bindings/dynamic"],
    "user_sources": ["../lib/bindings/dynamic/auto_dynamic.cpp"],
    "module_init": { "headers": ["auto_dynamic.h"],
                     "statements": ["bank::register_all()"] },
    "targets": ["python", "lua", "node"]
}
\end{lstlisting}

Six keys, every one about \emph{this} project. Read them against the diagram:
\field{user\_sources} is stage~1's output, and \field{module\_init} is the
\code{register\_all()} that publishes it.

\section{What it buys}

\begin{itemize}
  \item \textbf{One wrapper per language instead of one per class.} Ship the
    module once; any library that has been through the \lang{dynamic} backend
    becomes scriptable with no further code generation.
  \item \textbf{UI by query, not by codegen.} A property sheet is a walk over
    \code{fields()}: \code{doc} to tooltip, \code{range} to slider bounds,
    \code{combobox} and enumerators to a drop-down, \code{readonly} to a
    disabled row. The dispatch that a C++ property panel spends a file on is
    about ten lines of Python --- editable without recompiling anything.
  \item \textbf{Overloads come back} for the name-keyed targets. Lua, Node,
    Wasm, C\# and Java bind the first overload of a name and drop the siblings
    at generation time; through this module they do not, because
    \code{resolve()} scores the whole set at call time and can explain a
    failure.
  \item \textbf{Annotations are enforced once}, in the core, rather than
    re-implemented per backend --- a \code{range} violation comes back as the
    same message in every language.
\end{itemize}

That third point is worth dwelling on, because it is the sharpest evidence that
the IR is genuinely target-neutral. The same reflected overload set survives
into both projections; the expressiveness difference between them is
attributable entirely to \emph{when} names are resolved, not to information
lost on the way.

\begin{center}
\emph{The set of C++ members a target can express is not a property of the
target language. It is a property of when the target resolves names.}
\end{center}

\section{Running the pipeline twice}

\code{examples/plain-binding} runs it \textbf{once}: one library, one manifest,
four language bindings. That is the base case, and it is what most of this book
describes.

\code{examples/scriptable-model} runs stages 1--2 once per \emph{manifest}, and
it has two:

\begin{enumerate}
  \item \textbf{The user library to metadata.} The \lang{dynamic} backend turns
    \code{scene.h} into \code{auto\_dynamic.\{h,cpp\}} --- the same information
    the other backends spend on framework calls, emitted as plain data instead.
  \item \textbf{\code{<rosetta/script.h>} to bindings.} A second pipeline that
    binds rosetta's own reflection API, taking stage~1's output as its
    \field{user\_sources}.
\end{enumerate}

Practically: two manifests, two \code{rosetta\_gen --build} invocations, in that
order. Stage~2 must be re-run when stage~1's \emph{output location} changes ---
but not when your library's \emph{classes} change, which is the whole point.
